%! Author = layne
%! Date = 4/12/26

% Preamble
%  preamble
\documentclass[10pt]{article}

% dependencies
\usepackage[margin=0.6in]{geometry}
\usepackage{xcolor}
\usepackage[most]{tcolorbox}
\usepackage{enumitem}
\usepackage{multicol}
\usepackage{amsmath,amssymb}
\usepackage{array}
\usepackage{tabularx}

\tcbuselibrary{breakable}

% definitions
\definecolor{sky}{HTML}{EAF3FB}
\definecolor{navy}{HTML}{1F3A5F}
\definecolor{redbox}{HTML}{FCECEC}
\definecolor{greenbox}{HTML}{EDF8F0}
\definecolor{goldbox}{HTML}{FFF7E6}
\newtcolorbox{skillbox}[2][]{
    colback=#2,
    colframe=navy,
    boxrule=0.8pt,
    arc=2mm,
    left=2mm,right=2mm,top=1.5mm,bottom=1.5mm,
    fonttitle=\bfseries,
    title=#1,
    breakable
}

\setlist[itemize]{leftmargin=1.2em,itemsep=0.35em,topsep=0.35em}
\setlist[enumerate]{leftmargin=1.4em,itemsep=0.35em,topsep=0.35em}

% Centered column types
\newcolumntype{C}[1]{>{\centering\arraybackslash}p{#1}}
\newcolumntype{Y}{>{\centering\arraybackslash}X}

\newcommand{\CourseName}{Algebra 2}
\newcommand{\SchoolYear}{2026--2027}
\newcommand{\MeetingLength}{55 minutes}
\newcommand{\UnitNumberName}{Unit 1: Linear \& Absolute Value Functions}
\newcommand{\LessonNumberName}{Lesson 1: Linear Functions }

\begin{document}
    \begin{center}
        {\Large\bfseries \CourseName: \SchoolYear (\MeetingLength) \\
        {\small \bfseries \UnitNumberName \LessonNumberName } }
    \end{center}

    \begin{tcolorbox}[colback=sky,colframe=navy,boxrule=0.9pt,arc=2mm,left=2mm,right=2mm,top=1.5mm,bottom=1.5mm]
        \textbf{Primary objective:} Students will reactivate foundational algebra skills and connect prior knowledge to the study of functions in Algebra 2.
    \end{tcolorbox}

    \renewcommand{\arraystretch}{1.6}
    \begin{skillbox}[\centering Highest Priority Knowledge]{goldbox}
    \renewcommand{\arraystretch}{1.6}
    \begin{tabularx}{\linewidth}{C{0.32\linewidth} Y}
    \textbf{Skill} & \textbf{Why it matters} \\
    \hline
    Order of operations & Foundational fluency \\ \hline
    Integer operations & Cognitive load reduction \\ \hline
    Solving linear equations & Foundational fluency \\ \hline
    Simplifying expressions & Foundation fluency \\ \hline
    Operations with fractions & Rational expressions \\ \hline
    Factoring basics & Quadratics/Difference of squares \\ \hline
    Exponent rules & Polynomials/Exponentials/Logarithms \\ \hline
    Function notation & Math notation \\ \hline
    Graph interpretation & Parent functions and transformations \\ \hline
    Slope/rate of change & Functions/modeling \\ \hline
    Linear and Quadratic Regression & Functions/modeling \\
    \end{tabularx}
    \end{skillbox}




    \renewcommand{\arraystretch}{1.6}

    \begin{multicols}{2}
        \begin{skillbox}[1. Simplify \& Evaluate Expressions]{greenbox}
            \textbf{If there is one variable:}
            \begin{itemize}
                \item Use \textbf{$x$}.
                \item Type the expression as a function: $f(x)= [expression]$.
                \begin{itemize}
                    \item \textbf{Example:} find the value of $x^2-4$ when $x=7$.
                    \item In Desmos type in, $f(x)=x^2-4$
                \end{itemize}
                \item If a value is given, evaluate it with $f(value)$.
                \item If no value is given, graph the expression and compare it to the answer choices.
            \end{itemize}
            \textbf{If there is more than one variable:}
            \begin{itemize}
                \item Use variables such as \textbf{$a$, $b$, $c$}.
                \item Assign values like \textbf{$a=2$, $b=3$, $c=4$},
                but do \textbf{not} use $-1$, $0$, or $1$.
                \item Evaluate the original expression and then test answer choices until one matches.
            \end{itemize}

            \textbf{Useful inputs:}
            \begin{itemize}
                \item \textbf{$ f(x)=3x^2-2x+5 $}
                \item \textbf{$ f(2) $}
                \item $a=2$ \hspace{1em} $b=3$
            \end{itemize}

            \textbf{Why avoid $-1$, $0$, $1$?} Those values can make different expressions look the same by accident.
        \end{skillbox}

        \begin{skillbox}[2. Equations \& Inequalities]{redbox}
            \textbf{If there is one variable:}
            \begin{itemize}[noitemsep, topsep=0pt]
                \item Use \textbf{$x$}.
                \item Graph both sides as separate equations: $y=[left side]$ and $y=[right side]$.
                \item The \textbf{intersection} is the solution.
            \end{itemize}

            \textbf{If there are variables on both sides:}
            \begin{itemize}[noitemsep, topsep=0pt]
                \item Split the equation into two lines.
                \item Check what the graphs do:
                \begin{itemize}
                    \item intersect once $\rightarrow$ one solution
                    \item parallel $\rightarrow$ no solution
                    \item same line $\rightarrow$ infinite solutions
                \end{itemize}
            \end{itemize}

            \textbf{If an equation is set equal to 0:}
            \begin{itemize}
                \item Graph only the expression: $y=[expression]$.
                \item Look for \textbf{x-intercepts}.
                \item No x-intercept means \textbf{no real solution}.
            \end{itemize}

            \textbf{Inequalities:}
            \begin{itemize}
                \item Graph the inequality directly and use the shaded region.
                \item For systems of inequalities, the double-shaded region is the solution.
            \end{itemize}
        \end{skillbox}

        \begin{skillbox}[3. Regressions for Linear / Quadratic / Exponential Functions]{sky}
            \begin{itemize}
                \item {Use regression when you need the slope or equation of a line.}
                \item {Use regression when you need the equation of a quadratic/parabola}
                \item {Use regression when you need the equation of an exponential function}
            \end{itemize}
            \begin{itemize}
                \item Enter two points in a table (3 or more for quadratic/exponential).
                \item Run the appropriate regression using the regression utility
                \item Desmos will give you the equation for the regression.
            \end{itemize}

            \textbf{This is useful when:}
            \begin{itemize}
                \item you are given two or more points
                \item you can read two or more points from a graph
                \item the curve is written in one form and you want to use it in another.
            \end{itemize}

            \textbf{Strategy:}
            \begin{itemize}[noitemsep, topsep=2pt]
                \item Graph the line if needed.
                \item Pick two or more clear points on the line.
                \item Put them in a table.
                \item Use regression to get the equation.
            \end{itemize}
        \end{skillbox}

        \begin{skillbox}[4. Normal Distribution; Permutations \& Combinations]{sky}
            \begin{itemize}
                \item Univariate data: one variable
                \item Bivariate data: two variables
                \item Multivariate data: 3 or more variables
                \item \textbf{To determine what kind of data you have, ask yourself how many items you must return in an answer}
                \item For example, a study of the distribution of children's heights from ages 8-13, if you were asked this question about your child, you would give one answer -- the child's height.  Therefore, this is univariate data.
            \end{itemize}
            \textbf{Data Distributions}
            \begin{itemize}
                \item Normal:  mean=median=mode, bell-curve
                \item Positive/right skew:
                \item Negative/left skew:
            \end{itemize}

            \textbf{This is useful when:}
            \begin{itemize}
                \item you are given two or more points
            \end{itemize}

            \textbf{Strategy:}
            \begin{enumerate}
                \item Graph the line if needed.
            \end{enumerate}
        \end{skillbox}

        \begin{skillbox}[High-Impact Reminders]{goldbox}
            \begin{itemize}
                \item You \textbf{must} verify what you type into Desmos, and you must do it \textbf{every single time}.
                \item Use \textbf{$x$} for single-variable work.
                \item Split equations into two sides when Desmos does not show the answer clearly.
                \item If you cannot see a solution, zoom out or check x-intercepts.
                \item Pick easy-to-read points when using regression.
                \item When in doubt: \textbf{graph it, split it, or test values}.
            \end{itemize}
        \end{skillbox}

        \begin{skillbox}[Desmos Keyboard Shortcuts]{greenbox}
            \begin{itemize}
                \item sqrt = $\sqrt{}$
                \item cbrt = $\sqrt[3]{}$
                \item \texttt{\textasciicircum} exponent (shift-6)
                \item $f(x) = x^2 \{x \ge 0\}$ domain restriction
                \item pi = $\pi$
                \item gcf(num1, num2, num3, \ldots) = greatest common factor
                \item lcm(num1, num2, num3, \ldots) = lowest common multiple
                \item round(num1, \# places) = round number to \# places
            \end{itemize}
        \end{skillbox}
    \end{multicols}

\end{document}